\documentclass[10pt,a4paper]{amsart}
\usepackage[margin=19mm]{geometry}
\usepackage{amsmath,amssymb,mathtools}
\usepackage[scr=rsfs,cal=euler]{mathalfa}
\usepackage{xfrac}
\usepackage[unicode,hidelinks,bookmarks=false]{hyperref}
\newcommand{\ie}{i.e.\ }
\newcommand{\cf}{cf.\ }
\newcommand{\resp}{respectively\ }
\newcommand{\Subsection}{\subsection}
\providecommand{\nobreakdash}{\nobreak\hbox{-}}
\setlength{\emergencystretch}{2em}
\multlinegap0pt
\usepackage{amssymb}
\usepackage{xfrac}
\newtheorem{prop}{Proposition}
\title{Distance, Normals and Double Normals for real plane Curves with Singularities}
\author{Dirk Siersma}
\date{Source: arXiv:2606.08072v1 (CC BY 4.0)}
\begin{document}
\maketitle

\noindent\emph{Source and licence.} Distance, Normals and Double Normals for real plane Curves with Singularities. Version arXiv:2606.08072v1 (CC BY 4.0), distributed by arXiv under the Creative Commons Attribution 4.0 International licence, http://creativecommons.org/licenses/by/4.0/, as recorded in the arXiv licence metadata for this exact version and shown on the abstract page https://arxiv.org/abs/2606.08072v1. Full licence notice: \texttt{licenses/arxiv-2606.08072-CC-BY-4.0.txt}.\par\smallskip

\noindent\textit{Editorial scope.} Counted unit: the article's prop p:double-crit (source0.tex lines 675--684). The article's complete original proof of it is delivered with it (source0.tex lines 685--783, one \texttt{proof} environment of 99 lines). No mathematics is written, repaired or re-derived: the statement and the proof are the article's own lines, in the article's own notation, and no retained line is substituted or re-flowed.  Retained uncounted as the source-local context the proof needs: lines 167--169.  Nothing was omitted from any retained span, so the package carries no omission record.  No retained line contains a citation, so the wrapper carries no bibliography and no bibliography entry.  No retained line contains a figure environment, an image-inclusion command or drawing code, so no image is delivered and the package declares no asset dependency.  Editorial formatting only, in addition: an AMS \texttt{amsart} preamble that restates the article's own declarations and macros used by the retained lines, namely the environments \texttt{prop} on separate counters, together with the standard packages those lines need.  The retained environments print in this document as Proposition 1 in order of appearance, whereas the article prints the same items with the numbering of its own full document.  The complete native source of the article as distributed in the arXiv e-print for this exact version is bundled unchanged as \texttt{sources/202301/source0.tex}.
\begin{equation}\label{e:Puis}
x_1(t) = t ^p \; , \;  x_2(t) = c \,t^q + \sigma_{q+1}(t) 
\end{equation}

\begin{prop}\label{p:double-crit}
The top-critical points of $ \mathbb{B}$ are pairs $(a,b)$ of one of the following types:
\begin{itemize}
\item rr-normals (traditional double normals),  i.e. intervals $ab$ such that $ab$ is a normal in $a$ and $ba$ is  normal in $b$. Both $a$ and $b$ regular points of resp. $X_a$ and $X_b$,
\item rs-normal:   one end point is a geometric singular point of one curve;
 $ab$ is normal at the other point, which is assumed to be  regular,
 \item ss-normals: intervals between two geometric singular points of $X_a$ and $X_b$,
 \item the intersection points $X_a \cap X_b$.
\end{itemize}
\end{prop}

\begin{proof}
We work on the normalizations of the curves and use parametrizations of type (\ref{e:Puis}) at singular points. 
We choose the coordinates such that
$$ a(t) = (-1,0) +\varrho_{\alpha}(t^p, c_a t^q)   \;  \mbox{\rm with} \;  p <q  \;  \mbox{\rm and}\;  c_a = c_a[t] \; ; c_a[0] \ne 0  ;$$
$$ b(u) =\;  \; (1,0) + \varrho_{\beta} (u^r, c_b u^s)  \;  \mbox{\rm with} \;  r <s  \;  \mbox{\rm and}\;  c_b = c_b[u]\; ; c_b[0] \ne 0. $$
where $\varrho_{\alpha}$ is a rotation over $\alpha$. We have:
$$\mathbb{B} (t,u) =
-2 c_a c_b t^q u^s \cos (\alpha -\beta )-2 c_a t^q u^r \sin (\alpha -\beta )+c_a^2 t^{2 q}$$ $$-4 c_a \sin (\alpha ) t^q +2 c_b t^p u^s \sin (\alpha -\beta )+c_b^2 u^{2 s}+4 c_b \sin (\beta ) u^s$$ $$ - 2 t^p u^r \cos (\alpha -\beta )+t^{2 p} -4 \cos (\alpha ) t^p +u^{2 r}+4 \cos (\beta ) u^r+4.
$$
The dominant terms are:
\begin{equation}\label{dominant}
-4 \cos (\alpha ) t^p +4 \cos (\beta ) u^r.
 \end{equation}


There are two possible degenerations:
\begin{itemize}
\item $\cos (\alpha )= 0$, i.e.  $b$ lies on the normal to the curve $a(t)$ in $a$,
\item $\cos (\beta) = 0$, i.e. $a$ lies on the normal to the curve $b(u)$ in $b$.
\end{itemize}


Assume now: No degeneration:  the topological type of $\mathbb{B}$  at $(a,b)$ is determined by the dominant part.
We draw the following conclusions:
\begin{itemize}
\item If $p = odd$, or $r= odd$ then $\mathbb{B}$ is topologicaly regular (top-regular) at $(a,b)$,
\item If $p = even$ and $r = even$,
 then $\mathbb{B}$ behaves as  the Morse type $ -4 \cos (\alpha ) t^2 +4 \cos (\beta ) u^2$.
 It will have minimum, maximum or saddle, depending on the values of $\alpha$ and $\beta$.
\end{itemize}

{\bf Case rr:} Let us assume now that both curves are at least $C^1$-smooth. So we have $p=odd$ and $r=odd$.
If no degenerations: we have topological regularity. We are left with $\cos (\alpha) = 0 \;  \&  \;\cos (\beta) = 0$.
Which means that $ab$ is a double normal (in the usual sense). The dominant terms are now:
\begin{equation}\label{rr-condition}
 -4 c_a (-1)^k t^q +t^{2 p} +4 c_b (-1)^l u^s+u^{2 r}.
 \end{equation}
Note $  \alpha = \frac{\pi}{2} + k \pi$ and $\beta = \frac{\pi}{2} + l \pi$. 

We can distinguish between  the two independent cases:\\
$q < 2p \; ; \; q > 2p \; ; \; q = 2p$,\\
$s < 2r \; ; \; s > 2r \; ; \; s = 2r$.\\
This will result in: top-regular, maxima, saddle and minima for $\mathbb{B}$ . We leave this to the reader.
Note that in case $q=2p$ and $s=2r$ we get
\begin{equation}\label{rr-degenerate}
  -4 c_a (-1)^k t^{2p} +t^{2 p} +4 c_b (-1)^l u^{2r}+u^{2 r}.
\end{equation}
and the type depends on the positions of the two focal points. 
NB. This is the $C^2$-case and will be treated in terms of differential geometry and standard Morse theory below.

{\bf Case rs: } Lets us assume $a$ is a  smooth point and $b$ geometrically singular. We look again to the dominant terms (\ref{dominant}). Now $p = odd$ $r= even$: top-regular as long as $\cos(\alpha) \ne 0$.  Degenerate iff $ba$ is normal at $a$. In that case
the dominant and some nearest higher order terms are:
\begin{equation}\label{rs-condition}
  -4 c_a (-1)^k t^q +t^{2 p} +4 \cos (\beta ) u^r   +4 c_b \sin (\beta ) u^s+u^{2 r}
\end{equation}
Note $  \alpha = \frac{\pi}{2} + k \pi.$ 

As soon as $\cos (\beta) \ne 0$ The dominant terms are:
$$  -4 c_a (-1)^k t^q +t^{2 p} +4 \cos (\beta ) u^r.  $$
We can distinguish between
$q < 2p \; ; \; q > 2p \; ; \; q = 2p$:
\begin{itemize}
\item $q < 2p$: dominant: $-4 c_a (-1)^k  t^q +4  \cos(\beta ) u^r$ ; if $q$ is odd: top-regular, else minimum, maximum or saddle,
\item $q > 2p$: dominant: $ t^{2 p} +4 \cos (\beta ) u^r  $ \\
Since $r$ is even: minimum or saddle,
\item $q =2p$: dominant $  -4 c_a (-1)^k t^{2p} +t^{2 p} +4 \cos (\beta ) u^r  $.
 Since $r$ is  even: minimum, saddle or maximum depending on the  place of the focal points.
\end{itemize}




%%%%%%%%%%%%%%%%%%%%%%
{\bf Case ss:}   
 Lets us assume that  $a$  and $b$ are both geometrically singular. We look again to the dominant terms (\ref{dominant}). Now $p = even$ $r=even$. \textit{This is the generic case: The critical point is equivalent to a Morse critical point} as long as $\cos(\alpha) \ne 0$  and  $\cos(\beta) \ne 0$. Degenerate critical point of $\mathbb{B}$ iff $ba$ is normal at $a$ or $ab$ is normal at $b$. Assume the first, we can use again (\ref{rr-degenerate}):
dominant and nearest higher order terms are:
\begin{equation}
  -4 c_a (-1)^k t^q +t^{2 p} +4 \cos (\beta ) u^r +4 c_b \sin (\beta ) u^s+u^{2 r}.
\end{equation}
Note $  \alpha = \frac{\pi}{2} + k \pi.$ 
There is no difference with the  rs-case (\ref{rs-condition}).


If degenerate in both points $a$ and $b$ we get similar to the rr-case (\ref{rr-degenerate}):
$$  -4 c_a (-1)^k t^q +t^{2 p} +4 c_b (-1)^l u^s+u^{2 r}.$$
Note $  \alpha = \frac{\pi}{2} + k \pi$ and $\beta = \frac{\pi}{2} + l \pi$. 

We can distinguish between  the two independent cases:\\
$q < 2p \; ; \; q > 2p \; ; \; q = 2p$,\\
$s < 2r \; ; \; s > 2r \; ; \; s = 2r$.\\
This will result in: top-regular, maxima, saddle and minima . We leave this to the reader.
Note that in case $q=2p$ and $s=2r$ we get
\begin{equation}
 -4 c_a (-1)^k t^{2p} +t^{2 p} +4 c_b (-1)^l u^{2r}+u^{2 r}
\end{equation}
and the type depends on the positions of the two focal points.

{\bf Case intersection}: If the intersection of $X_a$ and $X_b$ is non-void then we have clearly an absolute minimum at each intersection point.
\end{proof}

\end{document}
